\chapter{Wykład XII}

\section{Ideał.}

\begin{context}

	Od teraz rozważając pierścień, rozważamy pierścień przemienny z $\mathbf{1}$.
\end{context}
Niech $R,S$ będą pierścieniami.
\begin{definition}[Jądro.]
	Niech $f: R \to S$ będzie homomorfizmem pierścieni. Definiujemy 
	\[
	\ker(f) := f^{-1}\left( \mathbf{0} \right) \text{ (jądro $f$ )}
	.\] 
\end{definition} 
\begin{fact}
	Mamy:
	\begin{itemize}
		\item $\ker(f) \le \left( R,+ \right) $.
		\item $\left(\forall r\in R \right)\left(\forall a \in \ker(f)  \right)\left( ra \in \ker(f)  \right)   $
	\end{itemize}
\end{fact}
\begin{proof}
	TODO
\end{proof} 
\begin{definition}[Ideał.]
	Niech $R$ jest pierścieniem oraz $\mathcal{I} \subseteq R$.

	$\mathcal{I}$ jest \textbf{ideałem} (oznaczamy $\mathcal{I} \unlhd R$ ), gdy:
	\begin{itemize}
		\item $\mathcal{I} \le \left( R,+ \right) $ (podgrupa grupy addytywnej $R$ ),
		\item $\left(\forall r\in R \right)\left(\forall a \in \mathcal{I} \right)\left( ra \in \mathcal{I} \right)   $
	\end{itemize}
\end{definition} 
\begin{remark}
	Ideały w $R$ to dokładnie $R$- podmoduły $R$.
\end{remark} 
\begin{eg}
	Niech $n\in \N$ wtedy $n\mathbb{Z} \unlhd \mathbb{Z} $ oraz $n\mathbb{Z} = \ker(r_{n}) $.

	Oczywiście $\ker(f) \unlhd R$ dla $f: R\to S$ będącego homomorfizmem pierścieni.
\end{eg} 
\begin{definition}[Warstwy.]
	Niech $\mathcal{I}\unlhd R$, wtedy przez $\mfaktor{R}{\mathcal{I}} $ określamy zbiór warstw addytywnych, dokładnie
	\[
	\mfaktor{R}{\mathcal{I}} := \{a+\mathcal{I}: a\in R\} 
	.\] 
\end{definition} 
\begin{theorem}[Warstwy pierścienia]
$\mfaktor{R}{\mathcal{I}} $ z działaniami 
\[
\left( a+ \mathcal{I} \right) + \left( b+\mathcal{I} \right) := \left( a+b \right) + \mathcal{I}
,\] 
oraz 
\[
	\left( a + \mathcal{I} \right) \cdot \left( b+\mathcal{I} \right) := \left( ab \right) + \mathcal{I}
,\] 
jest pierścieniem.

Natomiast 
 \begin{align*}
 	\Pi: R &\longrightarrow \mfaktor{R}{\mathcal{I}}  \\
 	r &\longmapsto r + \mathcal{I}
 .\end{align*}
jest homomorfizmem, oraz $\ker(\Pi) = \mathcal{I}$
\end{theorem} 
\begin{proof}
	TODO
\end{proof} 	
Niech $\mathcal{I} \unlhd R$ oraz $a,b \in R$.

Oznaczamy $a \equiv b \pmod{\mathcal{I}} $, gdy $a-b \in \mathcal{I} (\text{ czyli }a+\mathcal{I} = b+\mathcal{I})$.

Mówimy wtedy, że \textbf{$a$ przystaje do $b$ modulo $\mathcal{I}$}.

\begin{eg}
	$\mathcal{I} = n\mathbb{Z} \unlhd \mathbb{Z} $, $a,b \le \mathbb{Z} $.

	Wtedy
	\[
	\underbrace{a \equiv b \pmod{\mathcal{I}}}_{=n\mathbb{Z} } \iff \underbrace{a-b\in \mathcal{I}}_{=n\mathbb{Z} } \iff n\mid \left( a-b \right) \iff a \equiv b \pmod{n} 
.\] 
\end{eg} 
\section{Zasadnicze twierdzenie o homomorfiźmie pierścieni.}
\begin{definition}[Zasadnicze twierdzenie o homomorfiźmie pierścieni.]
	Niech $\varphi : R \to S $ będzie homomorfizmem pierścieni. Wtedy $\varphi \left[ R \right] \underset{\scriptscriptstyle\text{podp}}{\subseteq} S $ oraz istnieje dokładnie jeden izomorfizm pierścieni
	\[
	\psi : \mfaktor{R}{\ker(\varphi  ) \to \psi \left[ R \right] } 
	,\] 
	taki że następujący diagram komutuje (jest przemienny):

	\begin{center}
		\begin{tikzcd}[row sep= 30pt, column sep=30pt]
			R  \arrow[dr, "\Pi"] \arrow[rr, "\varphi"] & & \varphi \left[ R \right]     \\
		& \mfaktor{R}{\ker(\varphi) }  \arrow[ur, "\psi (\cong)", dashrightarrow] &  
		\end{tikzcd}
	\end{center}
\end{definition} 

\textbf{Przykłady}
\begin{enumerate}[label=\arabic*)]
	\item $r _{n} : \mathbb{Z} \to \mathbb{Z} _{n}, r _{n} \left[ \mathbb{Z}  \right] = \mathbb{Z} _{n} , \ker(r _{n} ) = n\mathbb{Z} $.

		Stąd $\mfaktor{\mathbb{Z} }{n\mathbb{Z} } \cong \mathbb{Z} _{n} $ 
	\item $R\left[ X \right] \overset{\text{ev}_{r}}{\longrightarrow} R, r\in R$.

		Mamy $\ker(\text{ev}_{r}) := \{f\in R\left[ X \right] : f\left( r \right) = \mathbf{0}\} $.

		Stąd $\mfaktor{R\left[ X \right] }{\{f\in R\left[ X \right] : f\left( r \right) = \mathbf{0}\} } \cong R$
\end{enumerate}

\begin{fact}

	Przekrój dowolnej rodziny ideałów w $R$ jest ideałem.
\end{fact}
\begin{proof}
	TODO
\end{proof} 
\begin{definition}
	Niech $R$ będzie pierścieniem, $A\subseteq R$ oraz $r_1,r_2,\ldots, r _{n} \in R$.
	\begin{enumerate}[label=\arabic*)]
		\item $\left( A \right) := \bigcap \{\mathcal{I}\unlhd R: A\subseteq \mathcal{I}\} $- najmniejszy ideał $R$ zawierający $A$.

			Inaczej ideał generowany przez  $A$ (sprawdź z rodziną Moore'a)
		\item $\left( r_1,r_2,\ldots,r _{n}  \right) := \left( \{r_1,r_2,\ldots,r_{n}\}  \right) $.
		\item $r_1R + r_2R + \ldots + r_{n}R := \{r_1a_1 + r_2a_2 + \ldots + r_{n}a_{n} : a_1,a_2,\ldots,a_{n} \in R\} $.

			$R$ jest to liniowa kombinacja elementów  $r_1,r_2,\ldots,r_{n}$.

			Jako zadanie pokazać, że 
			\[
				\left( r_1,r_2,\ldots, r_{n} \right) = r_1R + r_2R + \ldots + r_{n}R \text{ (ideał główny)} 
			.\] 
		\item Niech $\mathcal{I}\unlhd R$, jeśli istnieje $a\in \mathcal{I}$, takie że $\mathcal{I}=\left( a \right) $, to $\mathcal{I}$ nazywamy \textbf{ideałem głównym} .
		\item $\{\mathbf{0}\} $ oraz $R\unlhd R$ są \textbf{ideałami niewłaściwymi}, pozostałe ideały to \textbf{ideały właściwe}. 
	\end{enumerate}
\end{definition} 
\begin{fact}
	Dla $\mathcal{I} \unlhd R$ następujące warunki są równoważne:
	\begin{enumerate}[label=\roman*)]
		\item $\mathcal{I} = R$.
		\item $\mathbf{1} \in \mathcal{I}$.
		\item $\mathcal{I}\land R^{\ast} \neq \emptyset $
	\end{enumerate}
\end{fact}
\begin{proof}
	TODO
\end{proof} 
\begin{fact}

	$R$ jest ciałem $\iff R\neq \{\mathbf{0}\} $ i wszystkie ideały $R$ są niewłaściwe.
\end{fact}
\begin{proof}
	TODO
\end{proof} 
\begin{remark}
	Elementy $R \longrightarrow $ ideały $R$

	 $a \mapsto \left( a \right) $
\end{remark} 
\begin{fact}[Podzielność inkluzji ideałów]

	Niech $a,b \in R$. Wtedy
	\begin{enumerate}[label=\arabic*)]
		\item $a\mid b \iff \left( a \right) \supseteq \left( b \right) $.
		\item $a\sim b \left(\iff a\mid b \land b\mid a \right) \iff \left( a \right) = \left( b \right) $
	\end{enumerate}
\end{fact}
\begin{proof}
	TODO
\end{proof} 
\begin{definition}[PID]
	Pierścień $R$ nazywamy:
	\begin{enumerate}[label=\arabic*)]
		\item \textbf{Pierścieniem ideałów głównych}, gdy każdy ideał w $R$ jest główny.
		\item \textbf{Dziedziną ideałów głównych}(PID -Principal Ideal Domain?), gdy spełnia warunek $1)$ oraz jest dziedziną. 
	\end{enumerate}
\end{definition} 
\begin{eg}
	Mamy:
	\begin{enumerate}[label=\arabic*)]
		\item $K\text{ jest ciałem} \implies K \text{ jest PID}$, bo $\left(\forall \mathcal{I}\unlhd K \right)\left( \mathcal{I}= \underbrace{\{\mathbf{0}\} }_{= ( \mathbf{0} )} \text{ lub } \mathcal{I}= \underbrace{R}_{= (\mathbf{1})} \right)  $.
		\item $\mathbb{Z} \text{ jest PID}$.
	\end{enumerate}
\end{eg} 
\begin{proof}
	$\mathcal{I}\unlhd \mathbb{Z} $, załóżmy, że $\mathcal{I}\neq \{\mathbf{0}\} $. Weźmy $n\in \mathcal{I}$, takie że \[
	\left| n \right| = \min \{\left| x \right| : x\in \mathcal{I}\setminus \{\mathbf{0}\} \} 
	.\] 

	Wtedy 
	$\left(\forall b\in \mathcal{I} \right)\left(\exists q \in \mathbb{Z}  \right)\left( b = qn + r_{\left| n \right| }(b) \right)   .$ 

	Mamy $r_{\left| n \right| }(b) \in \{0,1,\ldots,\left| n \right| - 1\} $.

	Oraz $b, q_{n}\in \mathcal{I} \implies r_{\left| n \right| }(b) \in \mathcal{I}\implies r_{\left| n \right| }\left( b \right) = \mathbf{0}$.

	Zatem $b=qn$ i mamy $\mathcal{I} = \left( n \right) $
\end{proof} 

Powyższy dowód est istotny, bo użyjemy go do sformalizowania oewnego pojęcia.

W dowodzie mamy: $\{e\} \neq \mathcal{I} \unlhd \mathbb{Z} \implies \mathcal{I}= \left( n \right) $ dla pewnego $n\in \mathcal{I}$,takiego że $\left| n \right| = \min \{\left| x \right| : x\in \mathcal{I}\setminus \{\mathbf{0}\} \} .$ Istotne jest to, że w $\mathbb{Z}$ możemy dzielić z resztą względem $\left| \cdot  \right| $.

\begin{definition}
	Niech $R$ będzie pierścieniem.
	\begin{enumerate}[label=\arabic*)]
		\item Funkcja $v: R\setminus \{\mathbf{0}\} \to \N $ jest  \textbf{normą euklidesową} (funkcją euklidesową), gdy
			\[
			\left(\forall \underbrace{a}_{\neq \mathbf{0}},b\in R \right)\left(\exists \underbrace{q}_{\text{iloraz}}, \underbrace{r}_{\text{reszta}}\in R \right)\left( b = qa +r \land \left( v\left( r \right) < v\left( a \right) \lor r = 0 \right) \right)   
			.\] 
		\item $R$ jest \textbf{pierścieniem euklidesowym}, gdy istnieje norma euklidesowa na $R$.
	\end{enumerate}
\end{definition} 
\begin{eg}
	$\left| \cdot  \right| : \mathbb{Z} \setminus \{\mathbf{0}\} \to \N$ norma euklidesowa na $\mathbb{Z} $ bo
	\[
	\left(\forall \underbrace{a}_{\neq \mathbf{0}},b\in R \right)\left(\exists q,r = r_{\scriptscriptstyle \left| a \right| }\left( b \right)  \right)\left( b = qa+r \land \left( \left| r \right| = r < \left| a \right| \lor r = \mathbf{0} \right)  \right)   
	.\] 
\end{eg} 
\begin{theorem}
	Niech $v$ będzie normą euklidesową na $R$.

	Wtedy $\{\mathbf{0}\} \neq \mathcal{I}\unlhd R$ oraz $a\in \mathcal{I}\setminus \{\mathbf{0}\} $, taki że
	\[
	v\left( a \right) = \min \{v\left( b \right) : b \in \mathcal{I}\setminus \{\mathbf{0}\} \} 
	.\] 
	Wtedy $\mathcal{I} = \left( a \right) $.
\end{theorem} 
\begin{proof}
	TODO
\end{proof} 
\begin{corollary}

	$R$ jest euklidesowy $\implies$ $R$ jest pierścieniem ideałów głównych.
\end{corollary}
\begin{theorem}
	Niech $ K$ będzie ciałem. Wtedy
	\begin{enumerate}[label=\arabic*)]
		\item $K\left[ X \right] $ jest euklidesowy.
		\item $K\left[ X \right] $ jest PID.
	\end{enumerate}
\end{theorem} 
\begin{proof}
	TODO
\end{proof} 
\begin{remark}
	Analizując dowód, widzimy, że jeśli $a \in R\left[ X \right] $ jest taki, że wiodący współczynnik $a$ jest odwracalny, to możemy dzielić przez $a$ z resztą w sensie deg. Stąd mamy
\end{remark} 
\begin{theorem}[Bezout]
	Niech $R$ będzie pierścieniem, $\alpha \in R$, $g\in R\left[ X \right] $. Wtedy 
	 \[
	g\left( \alpha  \right) = \mathbf{0} \iff \left( X - \alpha  \right) \mid g \text{ w } R\left[ X \right] 
	.\] 
\end{theorem} 
\begin{proof}
	TODO
\end{proof} 
\begin{theorem}
	\begin{align*}
		\left| \cdot  \right| ^2: \mathbb{Z} \left[ i \right]  &\longrightarrow \N \\
		a+bi &\longmapsto a^2+b^2
	\end{align*}
	jest normą euklidesową na $\mathbb{Z} \left[ i \right] $.
\end{theorem} 
\begin{proof}
	TODO
\end{proof} 
\begin{corollary}

	$\mathbb{Z} \left[ i \right] $ jest PID.
\end{corollary}
\begin{remark}
	Rozważmy $\underbrace{d}_{\neq \mathbf{1}}\in \mathbb{Z} $ bezkwadratowa(???) oraz  \begin{align*}
		v: \mathbb{Q} \left( \sqrt{d}  \right)  &\longrightarrow \N \\
		n+m\sqrt{d}  &\longmapsto \left| n^2-m^2d \right|  
	.\end{align*}
	\begin{enumerate}[label=\arabic*)]
		\item Dla $d=-\mathbf{1}$, $v$ pokrywa się z $\left| \cdot  \right| ^2$ i jest normą euklidesową na $\mathbb{Z} \left[ \sqrt{d}  \right] = \mathbb{Z} \left( \mathbf{1} \right) $.
		\item $v$ jest normą euklidesową na pierścianiach
			\[
			\delta_{d}:= \begin{cases}
				\mathbb{Z} \left[ \frac{1+\sqrt{d} }{2} \right], &d \equiv 1 \pmod{4} \\
				\mathbb{Z} \left[ \sqrt{d}  \right], &d \not\equiv 1 \pmod{4} 
			\end{cases}
			.\] 
			Dla dokładnie 21 (bezkwadratowych) wartości \[d\in \mathbb{Z} : d\in \{-11,-7,-3,-1,2,3,5,6,\ldots,73\} .\]
		\item $v$ nie jest normą euklidesową na $\mathbb{Z} \left[ \sqrt{14}  \right] $, ale $\mathbb{Z} \left[ \sqrt{14}  \right] $ jest wciąż pierścieniem euklidesowym (względem innej normy).
		\item $\mathbb{Z} \left[ 19 \right] $ jest PID ale nie jest euklidesowa.
		\item $\mathbb{Z} \left[ \sqrt{-5}  \right] $ nie jest PID (udowodnimy później).
		\item To czy dany pierścień $\delta_{d}$ jest PID czy nie, jest istotne dla GFT (Great Fermat Theorem, LFO Last Fermat Theorem?
	\end{enumerate}
	Wiemy, że 
	\[
	K \text{ jest ciałem } \implies K\left[ X \right] \text{ euklidesowy } \implies K\left[ X \right] \text{ jest PID}
	.\] 
	Ale $\mathbb{Z} \left[ X \right] $ nie jest nawet PID (zadanie)
\end{remark} 
